\documentclass[11pt]{article}
\newcommand{\myname}{Jane Student}
\newcommand{\myemail}{jstudent}
\newcommand{\myhwtype}{Homework}
\newcommand{\myhwnum}{4}
\newcommand{\myclass}{CS 201}
\newcommand{\mylecture}{Lec}
\newcommand{\mysection}{A}

% Prefix for numbered questions.
\newcommand{\questiontype}{Question}

\usepackage{amsmath, amsfonts, amsthm, amssymb}
\usepackage{enumitem}
\usepackage{graphicx}
\usepackage{hyperref}
\usepackage{fancyhdr}
\usepackage[T1]{fontenc}
\usepackage{framed}

\usepackage{mathtools}
\DeclarePairedDelimiter{\ceil}{\lceil}{\rceil}
\DeclarePairedDelimiter{\floor}{\lfloor}{\rfloor}
\DeclarePairedDelimiter{\card}{\vert}{\vert}

\setlength{\parindent}{0pt}
\setlength{\parskip}{5pt plus 1pt}

\newcounter{questionCounter}
\newcounter{partCounter}[questionCounter]

% A named question, e.g. "Super Fancy Named Question", instead of a number.
\newenvironment{namedquestion}[1]{%
    \addtocounter{questionCounter}{1}%
    \setcounter{partCounter}{0}%
    \vspace{.2in}%
        \noindent{\bf #1}%
    \vspace{0.3em} \hrule \vspace{.1in}%
}{}

% An auto-numbered question.
\newenvironment{numedquestion}[0]{%
    \stepcounter{questionCounter}%
    \vspace{.2in}%
        \noindent{\bf \questiontype \; \arabic{questionCounter}. }%
    \vspace{0.3em} \hrule \vspace{.1in}%
}{}

\newenvironment{alphaparts}[0]{%
  \begin{enumerate}[label=\textbf{(\alph*)}]
}{\end{enumerate}}

\newenvironment{arabicparts}[0]{%
  \begin{enumerate}[label=\textbf{\arabic{questionCounter}.\arabic*})]
}{\end{enumerate}}

\newcommand{\answerbox}[1]{%
\begin{framed}
\vspace{#1}
\end{framed}}

\pagestyle{fancy}
\fancyhf{}
\fancyhead[L]{\textbf{\myclass\ \mylecture\mysection}}
\fancyhead[C]{\textbf{\myname\ (\myemail)}}
\fancyhead[R]{\textbf{\myhwtype\ \myhwnum}}
\renewcommand{\headrulewidth}{0.4pt}
\fancyfoot[C]{\thepage}

\begin{document}
\thispagestyle{plain}
\begin{center}
  {\Large \myclass{} \myhwtype{} \myhwnum} \\
  \myname{} (\myemail{}) \\
  \today
\end{center}

\begin{numedquestion}
  Prove that if $n$ is an odd integer, then $n^2$ is odd.

  \begin{framed}
  Write $n = 2k+1$ for some integer $k$. Then
  \[
    n^2 = (2k+1)^2 = 4k^2 + 4k + 1 = 2(2k^2 + 2k) + 1,
  \]
  which is of the form $2m+1$ with $m = 2k^2+2k$ an integer, so $n^2$ is odd.
  \end{framed}
\end{numedquestion}

\begin{numedquestion}
  Let $A$ and $B$ be sets. Show that $\card{A \cup B} = \card{A} + \card{B} -
  \card{A \cap B}$.

  \begin{framed}
  Every element of $A \cup B$ is counted once in $\card{A} + \card{B}$ unless
  it lies in both sets, in which case it is counted twice — once for $A$ and
  once for $B$. Subtracting $\card{A \cap B}$ removes exactly that double
  count, leaving each element of $A \cup B$ counted once.
  \end{framed}
\end{numedquestion}

\begin{namedquestion}{Recurrence Relation}
  Solve the recurrence $T(n) = 2T(n/2) + n$ with $T(1) = 1$.

  \begin{framed}
  By the Master Theorem with $a = 2$, $b = 2$, $f(n) = n$: since $f(n) =
  \Theta(n^{\log_b a}) = \Theta(n)$, this is case 2, so
  $T(n) = \Theta(n \log n)$.
  \end{framed}
\end{namedquestion}

\begin{numedquestion}
  Use the \texttt{arabicparts} environment for sub-parts that should include
  the question number.
  \begin{arabicparts}
    \item Show $\ceil{x} - \floor{x} \in \{0, 1\}$ for any real $x$.
    \item Give an $x$ for which this is $0$, and one for which it is $1$.
  \end{arabicparts}

  \begin{framed}
  \begin{arabicparts}
    \item If $x$ is an integer, $\ceil{x} = \floor{x} = x$, so the difference
    is $0$. Otherwise $\floor{x} < x < \ceil{x} = \floor{x} + 1$, so the
    difference is $1$.
    \item $x = 3$ gives $0$; $x = 3.5$ gives $1$.
  \end{arabicparts}
  \end{framed}
\end{numedquestion}

\begin{numedquestion}
  Use the \texttt{alphaparts} environment for lettered sub-parts, and prove
  the following by induction.
  \begin{alphaparts}
    \item State the claim: $\sum_{k=1}^{n} (2k - 1) = n^2$.
    \item Prove it.
  \end{alphaparts}

  \begin{framed}
  \begin{description}
    \item[Base case ($n=1$):] $\sum_{k=1}^{1}(2k-1) = 1 = 1^2$.
    \item[Inductive step:] Assume $\sum_{k=1}^{n}(2k-1) = n^2$. Then
    \[
      \sum_{k=1}^{n+1}(2k-1) = n^2 + (2(n+1)-1) = n^2 + 2n + 1 = (n+1)^2,
    \]
    completing the induction.
  \end{description}
  \end{framed}
\end{numedquestion}

% Use \setcounter{questionCounter}{<x>} to force the next question's number,
% e.g. when the written assignment's numbering does not start at 1.
\setcounter{questionCounter}{7}
\begin{numedquestion}
  This question is forced to start at number 8 — useful when a handout's
  numbering does not match the order questions appear in this document.

  \answerbox{2cm}
\end{numedquestion}

\end{document}
